\documentclass[10pt]{beamer}
\usetheme{Madrid}
\usepackage[T1]{fontenc}
\usepackage[utf8]{inputenc}
\usepackage[brazil]{babel}
\usepackage{lmodern}
\usepackage{hyperref}
\usepackage{listings}
\usepackage{xcolor}
\setbeamertemplate{navigation symbols}{}

\lstset{
  language=C++,
  basicstyle=\ttfamily\small,
  keywordstyle=\color{blue},\n  commentstyle=\color{gray},
  stringstyle=\color{red},
  breaklines=true,
  showstringspaces=false,
}

\title{Controle PID para Seguidor}
\subtitle{Curso de Microcontroladores}
\author{}
\date{}

\begin{document}

\begin{frame}
  \titlepage
\end{frame}

\begin{frame}{O problema sem PID}
  \begin{itemize}
    \item Abordagem simples: se desvia, vira para um lado.
    \item Resultado: robô oscila muito, zigzagueia.
    \item Pode sair completamente da linha em curvas rápidas.
    \item Sensível a ruído dos sensores.
  \end{itemize}
\end{frame}

\begin{frame}{O que é PID?}
  \begin{itemize}
    \item PID = Proporcional + Integral + Derivativo.
    \item É um controlador que faz ajustes suaves.
    \item Minimiza oscilação e erro de posição.
    \item Essencial em robótica.
  \end{itemize}
\end{frame}

\begin{frame}{P - Proporcional}
  \begin{itemize}
    \item Corrige proporcionalmente ao erro atual.
    \item Quanto maior o erro, maior a correção.
    \item Exemplo: se está muito desviado, vira bastante.
    \item Se está pouco desviado, vira pouco.
  \end{itemize}
\end{frame}

\begin{frame}{I - Integral}
  \begin{itemize}
    \item Acumula o erro ao longo do tempo.
    \item Corrige desvios sistemáticos persistentes.
    \item Se o robô fica sempre um pouco para a esquerda, integral corrige.
    \item Evita que o robô pare no meio da linha.
  \end{itemize}
\end{frame}

\begin{frame}{D - Derivativo}
  \begin{itemize}
    \item Detecta a velocidade de mudança do erro.
    \item Amortece oscilações rápidas.
    \item Se o erro está crescendo rápido, reduz a correção.
    \item Deixa o movimento mais suave.
  \end{itemize}
\end{frame}

\begin{frame}[fragile]{Fórmula básica do PID}
\begin{lstlisting}
// P: Proporcional
float P = erro * Kp;

// I: Integral (acumula ao longo do tempo)
integral += erro;
float I = integral * Ki;

// D: Derivativo (taxa de mudança)
float D = (erro - erro_anterior) * Kd;

// Sinal de controle final
float sinal = P + I + D;
erro_anterior = erro;
\end{lstlisting}
\end{frame}

\begin{frame}{Como ajustar PID}
  \begin{itemize}
    \item `Kp`: comece com valor pequeno, aumenta até oscilação.
    \item `Kd`: reduz oscilação, aumenta até ficar lento.
    \item `Ki`: última coisa a ajustar, pode desestabilizar.
    \item Processo: teste, observe, ajuste, repita.
  \end{itemize}
\end{frame}

\begin{frame}[fragile]{Exemplo simplificado em Arduino}
\begin{lstlisting}
float Kp = 0.1, Ki = 0.01, Kd = 0.05;
float integral = 0, erro_anterior = 0;

void loop() {
  int leitura = analogRead(SENSOR_CENTRO);
  float erro = leitura - 512;  // 512 é centro
  
  float P = erro * Kp;
  integral += erro;
  float I = integral * Ki;
  float D = (erro - erro_anterior) * Kd;
  
  float sinal = P + I + D;
  erro_anterior = erro;
  
  // Aplicar sinal aos motores
  int velocidade_esq = 150 - sinal;
  int velocidade_dir = 150 + sinal;
  
  analogWrite(MOTOR_ESQ, velocidade_esq);
  analogWrite(MOTOR_DIR, velocidade_dir);
}
\end{lstlisting}
\end{frame}

\begin{frame}{Por que PID?}
  \begin{itemize}
    \item Controle automático e robusto.
    \item Funciona mesmo com variações de velocidade.
    \item Fácil de entender e implementar.
    \item Padrão na indústria de robótica.
  \end{itemize}
\end{frame}

\end{document}
